% ECONOMICS_PROBLEM: {"id": "J61-503", "title": "Migration, selection and local adjustment", "jel_code": "J61", "source_id": "OP-0028"}
% FIRST_POSTED: 2026-10-09
% ATTEMPT_STATUS: partial
% ENTRY_KIND: research_attempt
\documentclass[11pt,a4paper]{article}
\usepackage[T1]{fontenc}
\usepackage[utf8]{inputenc}
\usepackage[margin=27mm]{geometry}
\usepackage{amsmath,amssymb,amsthm,mathtools}
\usepackage[hidelinks]{hyperref}
\setlength{\parindent}{0pt}
\setlength{\parskip}{5pt}
\newtheorem{theorem}{Theorem}[section]
\newtheorem{proposition}[theorem]{Proposition}
\newtheorem{lemma}[theorem]{Lemma}
\newcommand{\R}{\mathbb R}
\newcommand{\E}{\mathbb E}
\DeclareMathOperator{\Var}{Var}
\DeclareMathOperator{\Cov}{Cov}
\begin{document}
\section*{Migration: cohort outcomes and capital--housing adjustment}
\section{Scope and accounting}
This attempt derives an exact transition path in a specified receiving-economy benchmark, establishes sharp bounds when baseline natives are lost to follow-up, and separates earnings, housing and fiscal accounting. It does not estimate the effects of an actual admissions reform or solve an unrestricted dynamic spatial equilibrium. The baseline native population is fixed before the inflow. Tracking the current residents of a location is a different operation.

Consider $N>0$ native workers and a permanent inflow $M\ge0$ at time zero. All workers supply one unit of homogeneous labor; employment is full and nobody subsequently moves. This deliberately shuts down skill complementarity, selection and relocation. Competitive production is
$Y_t=A K_t^\alpha L^{1-\alpha}$, $0<\alpha<1$, $L=N+M$. Capital is financed externally and follows the announced adjustment rule
\[
 K_0=\bar kN,\qquad K_{t+1}=(1-\eta)K_t+\eta\bar k(N+M),\quad0<\eta\le1.
\]
Investment costs and returns belong to the outside suppliers; native welfare here includes labor earnings and specified housing positions. This rule is an explicit transition primitive, not a household Euler equation derived from saving behavior. In the no-inflow counterfactual, capital remains $\bar kN$. Let
$\bar w=(1-\alpha)A\bar k^\alpha$.

\section{An exact wage path and discounted loss bound}
\begin{proposition}
With $x=M/(N+M)$, the native wage ratio is
\[
 \frac{w_t}{\bar w}=\left[1-x(1-\eta)^t\right]^\alpha. \tag{1}
\]
For $0<\beta\le1$ and finite $H$, the discounted native earnings loss per worker obeys
\[
 \bar w\alpha x\sum_{t=0}^H[\beta(1-\eta)]^t
 \le \sum_{t=0}^H\beta^t(\bar w-w_t)
 \le \bar w x\sum_{t=0}^H[\beta(1-\eta)]^t. \tag{2}
\]
\end{proposition}
\begin{proof}
Solving the linear capital recursion gives
$K_t=\bar k[N+M-M(1-\eta)^t]$. Competitive marginal-product pricing yields (1). For $z\in[0,1]$ and $0<\alpha<1$, concavity gives $(1-z)^\alpha\le1-\alpha z$, while $(1-z)^\alpha\ge1-z$. Hence $\alpha z\le1-(1-z)^\alpha\le z$. Apply this with $z=x(1-\eta)^t$, multiply by $\bar w\beta^t$, and sum.
\end{proof}
If $\eta=1$, only the impact-period wage differs; subsequent capital has fully adjusted. If $\eta$ is small, the same eventual wage recovery masks a longer transition loss. For an infinite horizon, (2) remains finite whenever $\beta(1-\eta)<1$, including $\beta=1$ under positive adjustment. The formula supplies a testable path restriction, not a presumption that all immigration episodes obey it. In this homogeneous benchmark wages cannot exceed the initial level. Evidence of positive effects on a native skill group therefore requires an omitted margin such as complementarity, productivity, or heterogeneous tasks.

\section{Housing can prevent real-income recovery}
Suppose every household has Cobb--Douglas consumption preferences with housing expenditure share $0<\zeta<1$, the nonhousing price is one, and no household owns housing. The model is quasi-static in each period; there is no durable housing asset in the renter's budget. Let housing services supplied at rent $R$ be
$H_s(R)=H_0(R/R_0)^\varepsilon$, $\varepsilon\ge0$, where
$H_0=N\zeta\bar w/R_0$. Housing demand equals $L\zeta w_t/R_t$. Market clearing gives the unique positive rent
\[
 \frac{R_t}{R_0}=
 \left[\frac{N+M}{N}\frac{w_t}{\bar w}\right]^{1/(1+\varepsilon)}. \tag{3}
\]
The expenditure index for the stated preferences is a constant times $R_t^\zeta$, so native renter real-income ratios are
\[
 \frac{w_t/R_t^\zeta}{\bar w/R_0^\zeta}
 =\left(\frac{w_t}{\bar w}\right)^{1-\zeta/(1+\varepsilon)}
  \left(\frac N{N+M}\right)^{\zeta/(1+\varepsilon)}. \tag{4}
\]
Even as $w_t\to\bar w$, the limiting expression is below one when $M>0$ and finite housing elasticity. Unlike a fixed per-person housing requirement, this demand system clears even with fixed housing supply $\varepsilon=0$: people consume less housing per household at a higher rent. Equations (3)--(4) make that adjustment explicit.

If natives own housing, rental receipts must be added to their incomes before computing their consumption-equivalent welfare. Counting those receipts and also treating an asset-price capitalization of the same finite-horizon receipts as a separate gain would double count. Housing construction also uses resources; a total-surplus comparison must include them and the outside capital returns. The renter benchmark cannot be relabeled an all-native welfare result. Different tenure groups can have opposite incidence even when they face the same wage path.

\section{Tracking baseline natives: sharp missing-follow-up bounds}
For a policy regime $a$, let $F_i(a)$ indicate successful follow-up of a baseline native and suppose the declared outcome $Y_i(a)$, wherever the person lives, lies in $[\ell,u]$. Let $p_a=\Pr(F_i(a)=1)$ and $m_a=E[Y_i(a)\mid F_i(a)=1]$. A design identifying the regime's observed-data law gives
\[
 E[Y_i(a)]\in[p_am_a+(1-p_a)\ell,\quad p_am_a+(1-p_a)u]. \tag{5}
\]
To prove sharpness, leave the observed records fixed and give the missing natives any constant outcome between the two endpoints. Differences across regimes are bounded by subtracting opposite endpoints, since no cross-regime coupling is imposed. For a discounted earnings sum, use credible bounds on that whole sum or combine period bounds; missingness over time may make the latter conservative for joint paths.

A simple selection example shows the problem with current-resident averages. Two baseline natives initially earn one and three. Suppose the policy changes neither person's earnings, but the lower earner moves away. Local mean native earnings rise from two to three, although the baseline cohort's average treatment effect is zero. Adding a newly arriving high earner would further change the local mean without changing either incumbent's outcome. Following people through administrative identifiers addresses the population mismatch; it does not by itself establish the counterfactual assignment assumptions.

\section{Fiscal and innovation effects are separate unknowns}
Let a migrant's expected taxes minus explicitly allocated public services at duration $t$ be $b_t$. In a benchmark with no effects on native tax bases or congestion, the consolidated fiscal effect is $M\sum_{t=0}^H\beta^tb_t$. For a constant flow $b$, it is $Mb(1-\beta^{H+1})/(1-\beta)$ when $\beta<1$, and $Mb(H+1)$ when $\beta=1$. Intergovernmental transfers cancel in the consolidated account. If wage adjustment changes native tax receipts, that effect must be added, and the displayed simple formula no longer suffices. Fiscal receipts are not automatically net social benefits, since payments reduce somebody's consumption.

Nothing in (1)--(5) determines the innovation response. Adding $I_t=I_0+\gamma M$ with either sign of $\gamma$ leaves the wage and housing equations unchanged, so those equations alone cannot identify innovation. Similarly, observed short-run wage changes cannot distinguish low capital adjustment from a simultaneous productivity shock without extra measurements.

\section{Remaining gap}
The solved benchmark isolates capital speed and housing elasticity as distinct drivers, while the bounds reveal what missing movers can conceal. The catalogue additionally requires endogenous spatial choices, migrant selection, employment, fiscal interactions and validation on independent episodes. Estimating the benchmark on one inflow and reproducing that inflow is not such validation. The full multi-outcome policy vector remains unresolved.

\section{Completion Estimate}
35\%

This is a post hoc, heuristic estimate for the full catalogue target.
The attempt solves a homogeneous-labor capital and housing transition benchmark and bounds native outcomes missing after moves, with separate fiscal accounting. Full migration policy effects still lack endogenous migrant selection and location, employment and innovation responses, fiscal interactions, and validation on independent migration episodes.

\begin{thebibliography}{9}
\bibitem{card} D. Card (1990), \emph{The Impact of the Mariel Boatlift on the Miami Labor Market}, Industrial and Labor Relations Review 43, 245--257. Working-paper record: \url{https://www.nber.org/papers/w3069}.
\bibitem{op} G. I. P. Ottaviano and G. Peri (2012), \emph{Rethinking the Effect of Immigration on Wages}, Journal of the European Economic Association 10, 152--197. \url{https://doi.org/10.1111/j.1542-4774.2011.01052.x}. These motivate adjustment and substitution; the restricted transition equations above are not their empirical estimates.
\end{thebibliography}

\end{document}
